\section{来源审计：这堂课真正讲的是“生成分布如何被评估”}

Mark Chen 的讲座从 3-gram 一路走到 GPT-3、iGPT、DALL-E 与 Codex。表面上它像模型年表，真正的主线却是 evaluation interface 的不断变化：文本阶段比较 sample coherence 与人类识别；无监督学习阶段比较 transfer；图像阶段同时看 completion 与 representation；代码阶段则必须让生成结果通过 unit tests，并用 pass@k 评价多次采样。

原视频是全屏 slide talk，仓库没有独立 deck。本文从官方视频转场恢复 45 张 teaching slides；Stanford bumper、三个纯分隔页、致谢与重复 thank-you 被标为 optional。课堂发生于 2021 年，Codex、GPT-3 与 DALL-E 的产品/能力描述都按历史时点处理，不能倒写成 2026 年当前产品状态。

\begin{importantbox}{本讲的三个闭环}
\begin{enumerate}
  \item \textbf{Generation loop}：学习分布 $\rightarrow$ 采样 $\rightarrow$ 人类/任务评价 $\rightarrow$ 调整 scale 与 recipe；
  \item \textbf{Transfer loop}：unlabeled data $\rightarrow$ autoregressive objective $\rightarrow$ representation/behavior $\rightarrow$ downstream task；
  \item \textbf{Code loop}：problem prompt $\rightarrow$ sampled programs $\rightarrow$ unit tests $\rightarrow$ pass@k $\rightarrow$ temperature/model decision。
\end{enumerate}
\end{importantbox}

\begin{warningbox}{Sample quality、benchmark score 与 deployment reliability 不同}
一段文本像新闻，不证明它真实；一张图符合 prompt，不证明组合泛化；一个程序通过 benchmark tests，不证明它安全、可维护或满足隐藏需求。本文始终把“模型分布里存在正确样本”“当前 decoding 找到了样本”“部署系统可以信任输出”分开。
\end{warningbox}

\begin{knowledgebox}{术语消化：生成系统的四个对象}
\textbf{Model distribution} 是模型对候选输出的概率分布；\textbf{sample} 是从该分布得到的一次具体输出；\textbf{verifier} 用规则、tests 或人类评价判断候选；\textbf{selector} 在多个候选中选择最终输出。模型、采样、验证与选择可以分别改进，不能把最终成功全部归因于参数规模。
\end{knowledgebox}

\subsection{本章小结}

这不是单纯的 GPT 历史课，而是一堂生成建模评估课。随着输出从文本扩展到图像与代码，评价必须从表面相似度升级到任务证据和采样分布。

